\section{Related Work}
\label{sec:related}

\paragraph{Subtomogram averaging.}
Classic STA aligns and averages 3D subtomograms, compensating for the missing wedge, the region of Fourier space that the limited tilt range leaves unrecorded~\cite{forster2005av3,nicastro2006peet,castanodiez2012dynamo,hrabe2012pytom,wan2024stopgap}; RELION brought regularized-likelihood expectation-maximization over poses~\cite{scheres2012relion} to subtomograms~\cite{bharat2016relion}.
Later tools refine against the 2D tilt images.
EMAN2 corrects the CTF and refines the alignment of each particle in each tilt~\cite{chen2019eman2}; emClarity refines the tilt-series alignment with particles as fiducials and estimates per-tilt defocus~\cite{himes2018emclarity}; M~\cite{tegunov2021m} fits per-particle trajectories, tilt geometry and defocus to the tilt images; RELION-4/5~\cite{zivanov2022relion4,burt2024relion5} refines pseudo-subtomograms assembled from the tilt images and adds CTF refinement and Bayesian polishing~\cite{zivanov2019polishing}; nextPYP~\cite{liu2023nextpyp} refines constrained 2D particle projections.
Other methods improve the tomogram itself: NovaCTF~\cite{turonova2017novactf} corrects the defocus gradient through the sample depth during reconstruction, and gradient descent has recently been used to refine tilt-series alignment~\cite{chen2026tiltrefine}.
In all of these the map is a voxel grid rebuilt by back-projection, and pose search and each correction of geometry, optics and motion are separate steps, each followed by a new reconstruction.
% src: himes2018emclarity, chen2019eman2, tegunov2021m, zivanov2022relion4, burt2024relion5, liu2023nextpyp (all use voxel maps; refinement steps as described in each paper)

\paragraph{Neural representations.}
Neural networks are the other learned alternative to a voxel map.
CryoDRGN~\cite{zhong2021cryodrgn} and its cryo-ET versions~\cite{rangan2024cryodrgnet,powell2024tomodrgn} learn structural heterogeneity but take particle poses from a voxel pipeline; ICE-TIDE~\cite{debarnot2025icetide} and IsoNet~\cite{liu2022isonet} reconstruct or restore single tomograms without averaging.
None of them replaces the voxel map in pose search and refinement.

\paragraph{Gaussian models.}
Gaussian blobs and mixtures have long represented EM densities~\cite{marabini1998blobs,kawabata2008gmm}.
Joubert and Habeck~\cite{joubert2015gmm} infer a low-resolution initial model of Gaussian pseudo-atoms jointly with image orientations.
The e2gmm method~\cite{chen2021e2gmm,chen2023gmm} maps a mixture onto particle images in known orientations to model heterogeneity; a Gaussian mixture also refines the orientations of flexible regions~\cite{chen2024gmmrefine} and now serves as the alignment reference for focused refinement on 2D sub-tilt series in cryo-ET~\cite{chen2026gmmfocused}, while the map is still reconstructed from the particles.
3D Gaussian splatting~\cite{kerbl2023gaussian} has been adapted to tomography~\cite{zha2024r2gaussian} and to cryo-ET tomograms~\cite{zhang2026cryoetgs}.
CryoSplat~\cite{chen2026cryosplat} and GEM~\cite{qu2025gem} reconstruct homogeneous single-particle maps with Gaussians from known poses, and GaussianEM~\cite{he2025gaussianem} and CryoSPIRE~\cite{shekarforoush2025cryospire} model heterogeneity with deformable or hierarchical Gaussians.
Each of these methods uses Gaussians for one task.
STAGS instead fits a Gaussian mixture per half-set directly to the 2D tilt images and uses it as the map in every STA stage from the initial picks: global and local pose search, map fitting, and refinement of tilt geometry and particle motion.
